\section{Conclusion}\label{sec:conclusion}

{\color{red}
We introduced a diagnostic framework for comparing synthetic financial
return panels through four stylized facts.
Matched real--real and synthetic--real discrepancies provide a common
relative reference, while an outer ticker-corpus bootstrap and
independent inner-seed repeats distinguish corpus variation from
numerical variation.
The one-minute equity comparison identifies different strengths and
deficiencies across the four generators.
AIL has the smallest aggregate deviation and the highest scores on
the marginal distribution, aggregational Gaussianity, and conditional
tails, while MSV has the highest volatility-clustering score.
Only AIL's conditional-tail interval includes real--real gap parity;
this does not establish equivalence.
SBBTS improves on both controls for aggregational Gaussianity and
conditional tails but remains close to GBM on the other two properties.
The leading generator for each metric and the aggregate ordering
remain unchanged across 20 inner-resampling seeds.
The framework therefore complements aggregate assessments with an
interpretable account of where synthetic panels depart from observed
financial data.
The results constitute an in-sample empirical fidelity assessment of
fixed generator outputs over the full historical window, including the
SBBTS early-stopping overlap. The completed 19-specification analysis over
five paired seeds preserves the observed M1, M4 and aggregate orderings
within the tested ranges. Independent-period evaluation, repeated generator
fits and sensitivity to broader structural choices remain necessary extensions.
}
